\chapter{El Candado de Ayer: RSA y por qué las Computadoras Cuánticas lo Rompen}
\label{cap:rsa_y_shor}

\section{El mundo antes de la criptografía de retículos}

Antes de que existieran los retículos como herramienta criptográfica, el mundo digital se protegía con una familia de algoritmos basados en un tipo diferente de problema matemático: la \textbf{factorización de números enteros grandes}.

El representante más famoso de esta familia es \textbf{RSA}, creado en 1977 por Ron Rivest, Adi Shamir y Leonard Adleman (de ahí sus siglas). RSA fue el primer sistema de clave pública ampliamente adoptado y durante décadas protegió las transacciones bancarias, los correos electrónicos y los certificados de identidad digital de todo el mundo.

En este capítulo vamos a aprender cómo funciona RSA con números reales (¡solo necesitas saber dividir!), y luego descubriremos exactamente por qué el algoritmo cuántico de Shor es capaz de destruirlo en cuestión de segundos.

\section{El corazón de RSA: La multiplicación es fácil, la factorización no}

Empieza con un experimento mental:

\begin{center}
\begin{tabular}{ll}
\toprule
\textbf{Tarea} & \textbf{¿Cuánto tardas?} \\
\midrule
Multiplicar $13 \times 17 = ?$ & Menos de 5 segundos \\
Factorizar $221 = ? \times ?$ & Unos minutos (ensayo y error) \\
Multiplicar $4{,}294{,}967{,}311 \times 4{,}294{,}967{,}357 = ?$ & Segundos con calculadora \\
Factorizar $18{,}446{,}744{,}073{,}709{,}551{,}927$ & ¡Horas o días sin ayuda! \\
\bottomrule
\end{tabular}
\end{center}

\begin{metafora}[El producto de primos como una cerradura de combinación]
Imagina que tienes dos números primos $p$ y $q$ (números que solo son divisibles por sí mismos y por 1). Multiplicarlos es tan sencillo como abrir una puerta con su llave: instantáneo. Pero si solo te damos el producto $N = p \times q$ y te pedimos que encuentres de nuevo $p$ y $q$, tienes que probar cientos de miles de posibles divisores uno por uno.

Para números de 2048 bits (que tienen aproximadamente 617 dígitos decimales), el mejor algoritmo clásico de factorización tardaría más tiempo que la edad del universo en encontrar $p$ y $q$.
\end{metafora}

\begin{figure}[htbp]
\centering
\begin{tikzpicture}[scale=0.9]
  % Flecha facil izquierda a derecha
  \node[draw=PQCGreen, fill=BoxBgGreen, rounded corners=2mm, line width=1pt, minimum width=2.5cm, minimum height=1.2cm, align=center] (pq) at (0,0) {$p \times q$\\\small (dos primos)};
  
  \node[draw=PQCMainBlue, fill=BoxBgBlue, rounded corners=2mm, line width=1pt, minimum width=2.5cm, minimum height=1.2cm, align=center] (N) at (5,0) {$N$\\\small (producto público)};
  
  % Facil
  \draw[-{Stealth[length=3mm]}, line width=2pt, PQCGreen] (pq) -- (N) node[midway, above, font=\small\sffamily\bfseries, PQCGreen] {¡Fácil! (milisegundos)};
  
  % Dificil (vuelta)
  \draw[-{Stealth[length=3mm]}, line width=2pt, PQCCrimson, dashed] (N) to[bend left=40] node[midway, below, font=\small\sffamily\bfseries, PQCCrimson, align=center] {¡Difícil para clasicas!\\(siglos de cómputo)} (pq);
  
  % Flecha cuantica
  \draw[-{Stealth[length=3mm]}, line width=2pt, PQCAmber] (N) to[bend right=50] node[midway, above right, font=\small\sffamily\bfseries, PQCAmber, align=center] {Shor cuántico:\\segundos} (pq);
  
\end{tikzpicture}
\caption{La asimetría de la factorización: fácil hacia adelante, imposible hacia atrás para computadoras clásicas, pero trivial para Shor cuántico.}
\label{fig:factorizacion}
\end{figure}

\section{El funcionamiento de RSA paso a paso}

Veamos RSA con números pequeños para visualizar su mecánica real. En la práctica se usan números enormes (2048 bits), pero la estructura matemática es idéntica.

\subsection*{Generación de claves (lo que hace el receptor Roberto)}

\begin{ejemplonumerico}[RSA con números pequeños: $p=11$, $q=13$]
\begin{enumerate}
  \item \textbf{Elige dos primos:} $p = 11$ y $q = 13$.
  
  \item \textbf{Calcula el módulo público:} $N = p \times q = 11 \times 13 = 143$.
  
  \item \textbf{Calcula $\phi(N)$ (función de Euler):}
  \[
  \phi(N) = (p-1)(q-1) = 10 \times 12 = 120
  \]
  Esto cuenta cuántos números entre $1$ y $N$ no comparten factores con $N$.
  
  \item \textbf{Elige el exponente público $e$:} Escoge un número impar pequeño coprimo con $120$. Tomaremos $e = 7$.
  
  \item \textbf{Calcula el exponente privado $d$:} Busca $d$ tal que $e \cdot d \equiv 1 \pmod{\phi(N)}$, es decir, $7d \equiv 1 \pmod{120}$. La solución es $d = 103$ (ya que $7 \times 103 = 721 = 6 \times 120 + 1$).
  
  \item \textbf{Clave Pública:} $(e, N) = (7, 143)$.
  \item \textbf{Clave Privada:} $(d, N) = (103, 143)$.
\end{enumerate}
\end{ejemplonumerico}

\subsection*{Cifrado (lo que hace la emisora Alicia)}

Para cifrar un número $m$ (que representa nuestro mensaje), Alicia usa la clave pública de Roberto:
\[
c = m^e \pmod N = m^7 \pmod{143}
\]

\begin{ejemplonumerico}[Alicia cifra el mensaje $m=9$]
\begin{align*}
c &= 9^7 \pmod{143} \\
9^2 &= 81 \\
9^4 &= 81^2 = 6561 \equiv 6561 - 45 \times 143 = 6561 - 6435 = 126 \pmod{143} \\
9^7 &= 9^4 \cdot 9^2 \cdot 9 = 126 \times 81 \times 9 \pmod{143} \\
126 \times 81 &= 10206 \equiv 10206 - 71 \times 143 = 10206 - 10153 = 53 \pmod{143} \\
53 \times 9 &= 477 \equiv 477 - 3 \times 143 = 477 - 429 = 48 \pmod{143}
\end{align*}
Alicia envía $c = 48$.
\end{ejemplonumerico}

\subsection*{Descifrado (lo que hace Roberto con su clave privada)}

Roberto usa su exponente privado $d$ para recuperar el mensaje:
\[
m = c^d \pmod N = 48^{103} \pmod{143}
\]

Gracias al \textbf{Pequeño Teorema de Fermat} y la estructura del grupo multiplicativo $(\mathbb{Z}_N)^*$, este cálculo siempre devuelve el mensaje original $m = 9$.

\begin{intuicion}[¿Por qué funciona el descifrado?]
El matemático Euler demostró que para cualquier $m$ coprimo con $N = pq$:
\[
m^{\phi(N)} \equiv 1 \pmod N
\]
Por ello, $m^{e \cdot d} = m^{1 + k \cdot \phi(N)} = m \cdot (m^{\phi(N)})^k \equiv m \cdot 1^k = m \pmod N$.

La clave es que calcular $\phi(N) = (p-1)(q-1)$ requiere conocer $p$ y $q$. Un atacante que solo conoce $N$ no puede calcular $\phi(N)$ sin factorizar $N$. Y factorizar $N$ es el problema difícil.
\end{intuicion}

\section{La amenaza cuántica: El algoritmo de Shor}

En 1994, el matemático Peter Shor, trabajando en los Laboratorios Bell, publicó uno de los resultados más impactantes de la historia de la computación: un \textbf{algoritmo cuántico} capaz de factorizar cualquier número entero en tiempo polinomial.

Esto significa que una computadora cuántica suficientemente grande puede romper RSA-2048 en \textbf{horas o incluso minutos}, algo que a una computadora clásica le tomaría millones de años.

\subsection*{¿Qué hace especial a una computadora cuántica?}

\begin{metafora}[El laberinto de infinitos pasillos]
Imagina un laberinto gigante con millones de pasillos ramificados. Para un explorador clásico (una computadora normal), la única manera de encontrar la salida es recorrer los pasillos uno por uno o en pequeños grupos en paralelo.

Una computadora cuántica es como un explorador que puede desdoblarse en \textbf{millones de copias fantasmales simultáneas}, explorando todos los pasillos a la vez. Cuando uno de los fantasmas encuentra la salida, todos los demás colapsan hacia esa solución.
\end{metafora}

El principio físico detrás de esto es la \textbf{superposición cuántica}: a diferencia de los bits clásicos (que son 0 o 1), los \textbf{qubits} cuánticos pueden estar en un estado de superposición donde representan 0 y 1 \textit{al mismo tiempo}, hasta que se miden.

\begin{figure}[htbp]
\centering
\begin{tikzpicture}[scale=0.88]
  % Bit clasico
  \node[draw=gray!60, fill=gray!15, rounded corners=2mm, minimum width=2cm, minimum height=1.2cm, align=center] at (0,0) {\textbf{Bit clásico}\\$0$ \textit{ó} $1$};
  
  % Qubit
  \node[draw=PQCTeal, fill=BoxBgTeal, rounded corners=2mm, minimum width=2.5cm, minimum height=1.6cm, align=center] at (5,0) {\textbf{Qubit cuántico}\\$\alpha|0\rangle + \beta|1\rangle$\\\scriptsize (superposición)};
  
  % n clasicos
  \node[draw=gray!60, fill=gray!15, rounded corners=2mm, minimum width=3cm, minimum height=1.2cm, align=center] at (0,-2.5) {\textbf{$n$ bits clásicos}\\$2^n$ estados posibles,\\solo 1 activo a la vez};
  
  % n qubits
  \node[draw=PQCPurple, fill=BoxBgPurple, rounded corners=2mm, minimum width=3cm, minimum height=1.4cm, align=center] at (5,-2.5) {\textbf{$n$ qubits}\\$2^n$ estados activos\\\textbf{simultáneamente}};
  
  \draw[-{Stealth[length=2.5mm]}, line width=1.3pt, gray!60] (0,0) -- (0,-2.5);
  \draw[-{Stealth[length=2.5mm]}, line width=1.3pt, PQCTeal] (5,0) -- (5,-2.5);
  
\end{tikzpicture}
\caption{La diferencia fundamental: con $n=300$ qubits, una computadora cuántica puede explorar $2^{300}$ estados simultáneamente. Ese número supera la cantidad de átomos en el universo observable.}
\label{fig:clasico_vs_cuantico}
\end{figure}

\subsection*{El corazón de Shor: Detectar periodicidad oculta}

La genialidad del algoritmo de Shor radica en transformar el problema de factorizar $N$ en un problema de encontrar el \textbf{período de una función}, y luego usar el poder cuántico para encontrar ese período instantáneamente.

Aquí está la idea en etapas, sin fórmulas intimidantes:

\begin{figure}[htbp]
\centering
\begin{tikzpicture}[scale=0.9]
  % Columna de pasos
  \node[draw=PQCMainBlue, fill=BoxBgBlue, rounded corners=2mm, line width=1pt, align=center, minimum width=5cm] (s1) at (5, 7) {\textbf{Paso 1: Elegir $a$ al azar}\\\small Escoge un número $a$ con $1 < a < N$};
  
  \node[draw=PQCTeal, fill=BoxBgTeal, rounded corners=2mm, line width=1pt, align=center, minimum width=5cm] (s2) at (5, 5.4) {\textbf{Paso 2: Calcular la función periódica}\\\small $f(x) = a^x \pmod N$};
  
  \node[draw=PQCPurple, fill=BoxBgPurple, rounded corners=2mm, line width=1pt, align=center, minimum width=5.5cm] (s3) at (5, 3.8) {\textbf{Paso 3: Encontrar el período $r$}\\\small tal que $a^r \equiv 1 \pmod N$\\\small {\color{PQCGreen}$\Leftarrow$ ¡Aquí usa el poder cuántico!}};
  
  \node[draw=PQCGreen, fill=BoxBgGreen, rounded corners=2mm, line width=1pt, align=center, minimum width=5cm] (s4) at (5, 2.1) {\textbf{Paso 4: Calcular el factor}\\\small $\gcd(a^{r/2} - 1, N)$ = factor de $N$};
  
  \node[draw=PQCCrimson, fill=red!8, rounded corners=2mm, line width=1pt, align=center, minimum width=4cm] (s5) at (5, 0.5) {\textbf{¡RSA destruido!}\\\small Se conocen $p$ y $q$};
  
  \draw[-{Stealth[length=2.5mm]}, line width=1.3pt] (s1) -- (s2);
  \draw[-{Stealth[length=2.5mm]}, line width=1.3pt] (s2) -- (s3);
  \draw[-{Stealth[length=2.5mm]}, line width=1.3pt] (s3) -- (s4);
  \draw[-{Stealth[length=2.5mm]}, line width=1.3pt] (s4) -- (s5);
  
\end{tikzpicture}
\caption{Flujo simplificado del algoritmo de Shor. El paso crítico (detectar el período $r$) es exponencialmente difícil para computadoras clásicas pero se resuelve en tiempo polinomial con qubits cuánticos.}
\label{fig:shor_flujo}
\end{figure}

\subsection*{¿Por qué encontrar el período destruye RSA?}

Cuando encuentras el período $r$ de la función $a^x \pmod N$, ocurre algo algebraicamente hermoso. Si $r$ es par:
\[
a^r \equiv 1 \pmod N \quad \Rightarrow \quad (a^{r/2} - 1)(a^{r/2} + 1) \equiv 0 \pmod N
\]

Esto significa que $N$ divide el producto, pero no necesariamente cada factor individual. Al calcular el máximo común divisor:
\[
\gcd(a^{r/2} - 1, N) \quad \text{y} \quad \gcd(a^{r/2} + 1, N)
\]
con altísima probabilidad obtienes directamente $p$ y $q$. ¡El candado RSA queda abierto!

\begin{ejemplonumerico}[Shor con $N = 143$, $a = 2$]
Calculamos la secuencia $2^x \pmod{143}$:
\begin{center}
\begin{tabular}{cccccccc}
$x$: & 1 & 2 & 3 & 4 & $\cdots$ & 59 & 60 \\
$2^x$: & 2 & 4 & 8 & 16 & $\cdots$ & 33 & 66 \\
$x$: & 61 & 62 & $\cdots$ & \textbf{120} & & & \\
$2^x$: & 132 & 121 & $\cdots$ & \textbf{1} & & & \\
\end{tabular}
\end{center}
El período es $r = 120$. Como $r$ es par, calculamos:
\[
\gcd(2^{60} - 1, 143) = \gcd(\ldots, 143)
\]
Aplicando el algoritmo de Euclides se obtiene $13$. Y en efecto: $143 = 11 \times 13$.
\end{ejemplonumerico}

\section{¿Por qué una computadora cuántica encuentra el período tan rápido?}

Una computadora clásica necesitaría calcular $2^1, 2^2, 2^3, \dots$ hasta encontrar que $2^{120} \equiv 1$: hasta $120$ evaluaciones para este ejemplo pequeño, pero para un $N$ de 2048 bits serían potencialmente $2^{1024}$ evaluaciones.

La computadora cuántica aplica la \textbf{Transformada Cuántica de Fourier} (QFT), un equivalente cuántico de la Transformada Rápida de Fourier (FFT) que ya conoces de señales y audio:

\begin{figure}[htbp]
\centering
\begin{tikzpicture}[scale=0.92]
  % Funcion periodica (onda)
  \begin{scope}[shift={(0,0)}]
    \draw[line width=1.5pt, PQCTeal, domain=0:8, samples=200] plot (\x, {0.6*sin(\x*180/3.14*2/4)+1.5});
    \node[font=\small\sffamily, PQCTeal] at (4, 2.4) {$f(x) = a^x \pmod N$ (periódica con período $r$)};
    \draw[dashed, PQCAmber, line width=1pt] (0,0.3) -- (8,0.3);
    \draw[{Stealth[length=2mm]}-{Stealth[length=2mm]}, PQCAmber, line width=1pt] (0,0.3) -- (4,0.3) node[below, font=\tiny] {período $r$};
  \end{scope}
  
  % QFT
  \node[draw=PQCPurple, fill=BoxBgPurple, rounded corners=3mm, minimum width=2cm, minimum height=1cm, font=\small\sffamily\bfseries] (qft) at (9, 1.5) {QFT cuántica};
  
  % Espectro
  \begin{scope}[shift={(11,0)}]
    \foreach \x/\h in {1/0.2, 2/0.2, 3/2.2, 4/0.2, 5/0.2, 6/2.1, 7/0.2} {
      \fill[PQCGreen!80] ({\x*0.9-0.45}, 0) rectangle ({\x*0.9}, \h);
    }
    \node[font=\small\sffamily, PQCGreen] at (3.15, 2.7) {Picos = $N/r$};
  \end{scope}
  
  % Flecha QFT
  \draw[-{Stealth[length=2.5mm]}, line width=1.5pt, PQCPurple] (8.2, 1.5) -- (qft);
  \draw[-{Stealth[length=2.5mm]}, line width=1.5pt, PQCPurple] (qft) -- (10.1, 1.5);
  
\end{tikzpicture}
\caption{La QFT cuántica detecta el período $r$ de $f(x)$ en un solo paso: busca los picos de frecuencia en el espectro, que están separados exactamente por $N/r$.}
\label{fig:qft}
\end{figure}

La QFT coloca todos los $2^n$ posibles valores de $x$ en superposición simultánea, evalúa $f(x)$ en todos ellos a la vez, y luego extrae la información de periodicidad mediante interferencia cuántica. El período $r$ emerge de los picos del espectro.

Para un número de 2048 bits, la QFT necesita apenas del orden de $(2048)^2 \approx 4 \times 10^6$ operaciones cuánticas, en lugar de los $2^{1024}$ pasos de un algoritmo clásico.

\section{¿Y los retículos? La inmunidad a Shor}

La pregunta natural es: ¿por qué el algoritmo de Shor no rompe también a ML-KEM y ML-DSA?

\begin{intuicion}[La diferencia fundamental entre RSA y retículos]
\begin{itemize}
  \item \textbf{RSA} esconde el secreto en la \textbf{estructura periódica} del grupo multiplicativo $(\mathbb{Z}_N)^*$. La función $a^x \pmod N$ es perfectamente periódica, limpia, como una onda de radio pura. La QFT detecta esa onda en un instante.
  
  \item \textbf{Module-LWE} esconde el secreto en \textbf{datos ruidosos}: $\mathbf{b} = \mathbf{As} + \mathbf{e}$. El error $\mathbf{e}$ destruye completamente cualquier estructura periódica. No hay onda, solo caos aleatorio. La QFT no tiene nada que sintonizar.
\end{itemize}
\end{intuicion}

\begin{figure}[htbp]
\centering
\begin{tikzpicture}[scale=0.92]
  % RSA: onda limpia
  \begin{scope}[shift={(0,2.5)}]
    \node[font=\small\sffamily\bfseries, PQCCrimson] at (4, 2) {\textbf{RSA}: función periódica limpia $\Rightarrow$ Shor la detecta};
    \draw[line width=1.8pt, PQCCrimson, domain=0:8, samples=200] plot (\x, {0.7*sin(\x*180/3.14*2/2)});
    \draw[dashed, PQCAmber, line width=1pt] (0,-0.9) -- (8,-0.9);
    \draw[{Stealth[length=2mm]}-{Stealth[length=2mm]}, PQCAmber] (0,-0.9) -- (2,-0.9) node[midway, below, font=\tiny] {período};
  \end{scope}
  
  % Ruido LWE
  \begin{scope}[shift={(0,0)}]
    \node[font=\small\sffamily\bfseries, PQCGreen] at (4, 1.2) {\textbf{LWE}: datos ruidosos sin período $\Rightarrow$ Shor no tiene nada};
    \foreach \x in {0.3,0.6,...,7.8} {
      \pgfmathsetmacro\yr{(rand)*0.7}
      \draw[line width=1.5pt, PQCGreen!70] (\x, 0) -- (\x, \yr);
    }
  \end{scope}
  
\end{tikzpicture}
\caption{La razón geométrica de la inmunidad post-cuántica: RSA tiene una señal periódica limpia que Shor detecta; LWE tiene ruido puro sin patrón que ningún algoritmo cuántico puede explotar.}
\label{fig:rsa_vs_lwe_shor}
\end{figure}

Formalmente, los mejores algoritmos cuánticos conocidos para atacar problemas de retículos (como BKZ cuántico) siguen siendo \textbf{exponenciales} en la dimensión del retículo. Para dimensiones de $n = 768$ como las usadas en ML-KEM-768, la complejidad de romper el sistema se estima en $2^{178}$ operaciones cuánticas, lo cual está firmemente fuera del alcance de cualquier tecnología imaginable.

\section{El estado actual de las computadoras cuánticas}

En 2024, las computadoras cuánticas más avanzadas tienen entre 400 y 1000 qubits \textbf{físicos}, pero con tasas de error tan altas que en la práctica equivalen a solo unos pocos \textbf{qubits lógicos} estables. Para ejecutar el algoritmo de Shor sobre RSA-2048 se estima que se necesitan aproximadamente \textbf{4 millones de qubits lógicos} sin errores.

\begin{intuicion}[¿Por qué migrar ahora si los qubits no existen aún?]
El concepto de \textbf{``cosechar ahora, descifrar después''} (\textit{Harvest now, decrypt later}) es una amenaza real: adversarios sofisticados pueden estar capturando y almacenando hoy el tráfico cifrado con RSA, esperando el día en que tengan suficientes qubits para descifrarlo retroactivamente.

Datos cifrados hoy con RSA-2048 que deban mantenerse secretos por 20 o 30 años (como archivos médicos, estrategias de estado o claves de infraestructura crítica) están en riesgo si no se migran ahora a algoritmos post-cuánticos como ML-KEM y ML-DSA.
\end{intuicion}

Por eso el NIST publicó sus estándares post-cuánticos en 2024, y por eso el proyecto de grado de la UMSA implementa ML-KEM y ML-DSA en lugar de RSA para el sistema de votación electrónica.

\begin{resumencapitulo}
\begin{itemize}
  \item RSA se basa en la dificultad de \textbf{factorizar} el producto de dos números primos grandes.
  \item El proceso de cifrado/descifrado RSA usa potenciación modular: $c = m^e \pmod N$ y $m = c^d \pmod N$.
  \item El algoritmo de Shor transforma la factorización en detección de periodicidad, usando la QFT cuántica para encontrar el período $r$ de $a^x \pmod N$ en tiempo polinomial.
  \item LWE/Module-LWE no tiene estructura periódica (el ruido $\mathbf{e}$ la destruye), por lo que Shor no puede aplicarse.
  \item La migración a criptografía post-cuántica es urgente hoy debido al riesgo de ``cosechar ahora, descifrar después''.
\end{itemize}
\end{resumencapitulo}
